\vspace{0.14in}
\begin{abstract}
This case study presents a rigorous empirical and theoretical comparison between variational quantum classifiers (QNN/VQC) with data re-uploading and classical multilayer perceptrons (MLPs). We evaluate models across four controlled benchmarks: an overlapping botanical classification task (\emph{Iris versicolor} vs. \emph{Iris virginica}), the nonlinear Two Moons distribution, a rotationally symmetric Concentric Circles task, and physical two-qubit time-dependent Hamiltonian dynamics from \emph{QDataSet} (Perrier et al., 2022). The experiment employs ten paired stratified train/test splits, capacity-matched compact classical networks as parameter-controlled baselines (9 parameters for 2D inputs, 13 parameters for 4D inputs), and wider MLPs (193 and 225 parameters) as expressive references. We isolate architectural effects via circuit depth scaling (depths 1--3), an entanglement ablation removing CNOT gates, and multinomial finite-shot inference sampling (256 and 1,024 shots). Executed with PennyLane \texttt{default.qubit} (Torch interface, simulator backpropagation) and verified against an exact PyTorch statevector reference, the study demonstrates that classical MLPs consistently outperform or match QNNs on generic classical tabular datasets: $0.956\pm0.030$ vs. $0.912\pm0.037$ on Iris ($p=0.0032$) and $0.938\pm0.024$ vs. $0.763\pm0.030$ on Two Moons ($p < 10^{-5}$). However, on Concentric Circles, trigonometric rotation encoding imparts a decisive inductive bias: depth-1 ($0.980\pm0.017$) and depth-3 QNNs ($0.977\pm0.023$) dramatically outperform the capacity-matched compact MLP ($0.584\pm0.109$, $p=0.0033$). On physical QDataSet Hamiltonian dynamics, the 13-parameter depth-2 QNN ($0.742\pm0.047$) reliably captures non-local interaction phases alongside the 13-parameter compact MLP ($0.818\pm0.082$), with depth-3 QNN scaling to $0.765\pm0.038$. We formulate the theoretical principles governing why classical NNs dominate classical learning (gradient dynamics, non-unitary activations, unconstrained parameter scaling) while detailing the four precise problem regimes where quantum models hold proven or plausible advantage (quantum-native physical dynamics, discrete-logarithm kernels, high-order parity, and classically hard quantum kernels). Turnkey reproducibility is provided via an automated cross-platform pipeline.
\end{abstract}
\vspace{0.04in}
\noindent\textbf{Keywords:} quantum machine learning, variational quantum circuit, QNN, inductive bias, classical neural network, QDataSet, PennyLane, PyTorch, benchmarking.

\vspace{0.08in}
\begin{center}
\begin{tikzpicture}[node distance=0.34cm and 0.34cm, font=\small]
  \node[arch,fill=qgold,text width=2.45cm] (data) {Benchmark data\\Iris / Moons / Circles\\QDataSet Dynamics};
  \node[arch,fill=qblue,right=0.68cm of data,text width=2.45cm] (mlp) {Classical path\\Standardize features\\PyTorch MLP};
  \node[arch,fill=qteal,below=0.55cm of mlp,text width=2.45cm] (qnn) {Quantum path\\Angle encoding\\2--4 qubit VQC};
  \node[arch,fill=qcoral,right=0.65cm of mlp,text width=2.15cm] (metrics) {Shared held-out test set\\Accuracy, F1, AUC\\and runtime};
  \draw[arrow] (data.east) -- ++(0.22,0) |- (mlp.west);
  \draw[arrow] (data.east) -- ++(0.22,0) |- (qnn.west);
  \draw[arrow] (mlp.east) -- (metrics.west);
  \draw[arrow] (qnn.east) -| ([yshift=-0.11cm]metrics.west);
\end{tikzpicture}
\end{center}
\vfill
\noindent\textbf{Scope note.} The primary comparison evaluates simulated QNN circuits vs classical MLPs on classical and quantum-simulated data to analyze inductive bias, trainability, and representation limits. It does not claim physical quantum hardware speedup.

